\section{Constexpr Lexing and Parsing}

\subsection{Zero-Copy Tokenizer}
The \texttt{ctpb::lexer} processes raw IDL text during constant evaluation. Unlike runtime tokenizers that allocate storage for lexemes, every identifier, keyword, and literal is retained as a \texttt{std::string\_view} referencing the embedded source, so tokenization performs zero dynamic allocations. The lexer recognizes the proto3 keyword set (\texttt{syntax}, \texttt{package}, \texttt{import}, \texttt{message}, \texttt{enum}, \texttt{oneof}, \texttt{map}, \texttt{reserved}), integer and floating-point literals including hexadecimal, single- and multi-line comments, and punctuation, in $\mathcal{O}(N)$ time.

\subsection{Recursive-Descent Parser}
The proto3 grammar is formalized into approximately fifty LL(1) productions, and the parser builds the structural AST in a single pass without backtracking. Message parsing dispatches on the leading token of each body element, recursing for nested messages and delegating to dedicated productions for enums, oneofs, maps, and reserved ranges; anything else is parsed as a field declaration.

\subsection{Compile-Time Diagnostics}
When a syntax violation, invalid field number, or reserved-field conflict is detected during \texttt{consteval} evaluation, \texttt{ctpb} raises an exception carrying the source line and column. Under ISO C++ constant-evaluation rules, throwing inside a \texttt{consteval} context terminates compilation, so the schema error surfaces as a localized compiler diagnostic rather than a runtime failure.
